\documentclass[11pt,a4paper]{amsart}

\usepackage[T1]{fontenc}
\usepackage{amsmath,amssymb,amsthm,mathtools}
\usepackage{enumitem}
\usepackage{microtype}
\usepackage{etoolbox}
\usepackage[margin=2.7cm]{geometry}
\usepackage{xcolor}
\definecolor{linkcol}{rgb}{0.10,0.20,0.50}
\usepackage[colorlinks=true,linkcolor=linkcol,citecolor=linkcol,urlcolor=linkcol]{hyperref}
\makeatletter
\g@addto@macro\UrlBreaks{\do\_\do\.}
\makeatother

\theoremstyle{plain}
\newtheorem{theorem}{Theorem}[section]
\newtheorem{proposition}[theorem]{Proposition}
\newtheorem{lemma}[theorem]{Lemma}
\newtheorem{corollary}[theorem]{Corollary}
\theoremstyle{definition}
\newtheorem{definition}[theorem]{Definition}
\newtheorem{example}[theorem]{Example}
\newtheorem{remark}[theorem]{Remark}

% Keep consecutive theorem/proof blocks visually distinct in the compact amsart layout.
\AfterEndEnvironment{theorem}{\par\smallskip}
\AfterEndEnvironment{proposition}{\par\smallskip}
\AfterEndEnvironment{lemma}{\par\smallskip}
\AfterEndEnvironment{corollary}{\par\smallskip}
\AfterEndEnvironment{proof}{\par\smallskip}

\newcommand{\RR}{\mathbb R}
\newcommand{\CC}{\mathbb C}
\newcommand{\NN}{\mathbb N}
\newcommand{\one}{\mathbf 1}
\newcommand{\rank}{\operatorname{rank}}
\newcommand{\corank}{\operatorname{corank}}
\newcommand{\indm}{\operatorname{ind}_{-}}
\newcommand{\Spec}{\operatorname{Spec}}
\newcommand{\diag}{\operatorname{diag}}
\newcommand{\spanr}{\operatorname{span}_{\RR}}
\newcommand{\Res}{\mathcal R}
\newcommand{\Mcal}{\mathcal M}
\newcommand{\Hcal}{\mathcal H}
\newcommand{\Tcal}{\mathcal T}
\newcommand{\Ecal}{\mathcal E}
\emergencystretch=1.5em

\begin{document}

\title[The Euler--capacity law]{The Euler--capacity law for wild character varieties}

\author{M.~Haufschild}
\address{Independent researcher}
\email{haufschildmichael@gmail.com}
\date{3 August 2026}

\subjclass[2020]{Primary 34M40, 14D24; Secondary 15A18, 32G34, 14H60, 16G20}
\keywords{irregular Riemann--Hilbert correspondence, wild nonabelian Hodge theory,
Stokes data, quasi-Hamiltonian fusion, rigidity index, fission tree, generalized Cartan matrix,
inertia, effective resistance, principal conormal, Rees degeneration}

\begin{abstract}
Let $X$ be a compact curve of genus $g$ with a nonempty set $D$ of $m$ marked points carrying
type-$A$ irregular data for general linear groups.  The dimension of the corresponding wild character
variety determines a quadratic index $Q_{X,D}$ on the real marked multiplicity lattice.  We prove
the fusion identity
\[
 Q_{X,D}=\sum_{a\in D}Q_a+2(1-g-m)M^2,
\]
where $Q_a$ is the one-puncture index and $M$ is the common total rank.  Each local index has a
canonical capacity completion $Q_a=P_a(z_a)+\Omega_aM^2$, with $P_a$ positive definite.  Hence
\[
 \boxed{\quad Q_{X,D}=\sum_{a\in D}P_a(z_a)
       +\left(\chi(X)-\sum_{a\in D}\delta_a\right)M^2,\qquad
       \delta_a=2-\Omega_a.\quad}
\]
We call this the Euler--capacity law: Euler topology minus local capacity defects.  It is a
quadratic refinement of scalar rigidity-index and
Euler-characteristic formulae: topology and all punctures couple only through one physical-rank
coordinate, while every centered Stokes direction has strictly positive energy.  It follows that
the negative index and nullity are at most one, in every genus, and that the sign of the displayed
coefficient gives the finite/affine/Lorentz trichotomy.  The coefficient is intrinsically the
minimum of $Q_{X,D}$ at unit common rank.

The local theorem making the formula effective is a Stokes contact-tree factorization
\[
 C=I+\sum_e\ell_eu_eu_e^T-(H-1)rr^T.
\]
Its positive squares are principal conormals of fission clusters.  The reciprocal local capacity
is the effective resistance of the contact tree; corona pruning is exact series--parallel
reduction; type-$A$ legs only change terminal resistances; and Rees scaling preserves the unique
rank defect.  In genus zero with one puncture the Euler--capacity law recovers the complete
Stokes--Cartan theorem, including the determinant wall and the graph obstruction
$\lambda_2(B)\leq2$.  We state exact hypotheses and falsifiers: arbitrary covectors need not be
rank-like, and no categorical equality with vanishing-cycle multiplicities is asserted.
\end{abstract}

\maketitle

\section{Introduction}

Wild nonabelian Hodge theory relates meromorphic connections, irregular Higgs bundles, and Stokes
local systems on curves \cite{BiquardBoalch2004}.  Their Betti moduli spaces are assembled by
quasi-Hamiltonian fusion from local Stokes spaces and surface handles
\cite{Boalch2014Geometry,BoalchYamakawa2015Twisted}.  In genus zero, marked formal data also determine a finite diagram with
a symmetric generalized Cartan matrix, whose quadratic form computes the dimension of the wild
character variety \cite{BoalchYamakawa2020,Doucot2021Diagrams}.  The diagram is invariant under
symplectic automorphisms of the Weyl algebra, including Fourier--Laplace transformation
\cite{Doucot2021Diagrams}.

There is a scalar topological precedent.  Under multiplicity-free and spectral regularity
hypotheses, the rigidity index of an irregular connection equals the Euler characteristic of the
normalization of its irregular spectral curve \cite{Hiroe2019Rigidity}.  Locally, irregularity,
intersection numbers, Milnor numbers, Fourier transform, and blow-up are linked by a plane-curve
dictionary \cite{Hiroe2017Ramified}.  These results suggest that the scalar index is the shadow of
a polarized law.  They do not determine the inertia of the full marked lattice, nor identify which
deformation directions can carry negative energy.

Our main result supplies that missing polarization.  Let $D=\{p_1,\dots,p_m\}$ be nonempty, and
at each $p_a$ fix a marked type-$A$ irregular datum of total rank $M$.  Denote its real marked
multiplicity space by $V_a$, its rank functional by $R_a$, and its one-puncture index by $Q_a$.
The global multiplicity space is the rank fibre product
\begin{equation}\label{eq:intro-fibre-product}
 V_{X,D}=\{(x_a)\in\textstyle\bigoplus_aV_a:R_1(x_1)=\cdots=R_m(x_m)\},
 \qquad M=R_a(x_a).
\end{equation}
The first universal identity is
\begin{equation}\label{eq:intro-fusion}
 \boxed{\quad Q_{X,D}(x)=\sum_{a=1}^mQ_a(x_a)+2(1-g-m)M^2.\quad}
\end{equation}
It follows by comparing the exact dimensions of the global quasi-Hamiltonian reduction and the
$m$ one-puncture reductions.  No stability inequality or asymptotic approximation enters it.

The second input is local and constructive.  Fission trees record how the formal structure group
breaks at an irregular pole and classify admissible deformation data
\cite{Boalch2008Kac,BoalchDoucotRembado2025,Boalch2025Counting}.  The resulting one-puncture graphs
range beyond finite and affine Dynkin diagrams; hyperbolic examples occur
\cite{Boalch2012Simply,Boalch2025Counting}.  Ramification allows loops or signed entries, and some
twisted irregular classes yield honest graphs that are not fission graphs
\cite{Doucot2025Graphs}.  Their universal local law has the short form
\begin{equation}\label{eq:intro-law}
 \boxed{\quad C=P-\lambda rr^T,\qquad P>0.\quad}
\end{equation}
The positive part $P$ is not an unspecified matrix.  It is the sum of a positive diagonal and one
rank-one Gram term for every fission cut.  The only subtractive term is the square of total
ramified rank.  Formula \eqref{eq:intro-law} immediately suggests negative index at most one, but
its force is greater: a single intrinsic scalar $\Omega$ detects all three possible signatures,
the affine kernel, the determinant, and the exact transition wall.

For the core matrix $B=(B_{ij})$ and ramification orders $r_i$, introduce the complementary
contact coordinates
\begin{equation}\label{eq:intro-K}
 K_{ij}=1+\frac{B_{ij}}{r_ir_j}\quad(i\ne j),\qquad
 K_{ii}=1+\frac{B_{ii}-1}{r_i^2}.
\end{equation}
The fission-contact hierarchy realizes $K_{ij}$ as the height of the meet of leaves $i$ and $j$.
Extend the root to any height $H>\max(1,K_{ij})$.  If $A_e$ is the descendant leaf set of an edge,
$\ell_e>0$ its contact length, and $u_e=(r_i\one_{A_e}(i))_i$, our first main theorem is
\begin{equation}\label{eq:intro-factor}
 C=I+\sum_e\ell_eu_eu_e^T-(H-1)rr^T.
\end{equation}
The extension to height $H$ is auxiliary.  Adding a root segment of length $t$ adds $t rr^T$ to
the positive Gram part and simultaneously replaces $H-1$ by $H-1+t$, leaving $C$ unchanged.

Set
\begin{equation}\label{eq:intro-capacity}
 P=I+\sum_e\ell_eu_eu_e^T,\qquad
 c=r^TP^{-1}r,\qquad
 \Omega=c^{-1}-(H-1).
\end{equation}
Then $c>0$, and every $x$ splits uniquely into a vector $z$ of zero total ramified rank and a
multiple of $P^{-1}r$.  In these coordinates,
\begin{equation}\label{eq:intro-centered}
 x^TCx=z^TPz+(r^Tx)^2\Omega.
\end{equation}
Consequently $C$ is positive definite for $\Omega>0$, positive semidefinite with a one-dimensional
kernel for $\Omega=0$, and nonsingular with exactly one negative eigenvalue for $\Omega<0$.
Neither multiple affine directions nor coexistence of affine and Lorentz directions is possible.
Throughout, ``Lorentzian'' means signature $(N-1,1)$ for a nonsingular real symmetric form on an
$N$-dimensional space; no stronger condition on all proper principal submatrices is built into the
term.

Applying this completion at every puncture in \eqref{eq:intro-fusion} yields the headline law
\begin{equation}\label{eq:intro-euler-capacity}
 \boxed{\quad Q_{X,D}(x)=\sum_{a=1}^mP_a(z_a)+\Omega_{X,D}M^2,
 \qquad
 \Omega_{X,D}=\chi(X)-\sum_{a=1}^m\delta_a,qquad
 \delta_a=2-\Omega_a.\quad}
\end{equation}
This is the one-port index theorem: $Q_{X,D}$ is positive definite on the hyperplane $M=0$,
so its entire nonpositive inertia is confined to the one-dimensional quotient detected by physical
rank.  Thus $\indm Q_{X,D}+\dim\ker Q_{X,D}\le1$, and the sign of $\Omega_{X,D}$ gives the exact
trichotomy.  Each handle subtracts $2$ from the same scalar channel.  Fusion adds inverse local
capacities; no new spatial defect is created.

There are five reasons the factorization is geometric rather than merely a rank-one perturbation
observation.
\begin{enumerate}[label=(\roman*),leftmargin=2.2em]
\item Each positive summand is the square of the principal conormal
$\nu_e(x)=\sum_{i\in A_e}r_ix_i$ of a canonical cluster hyperplane.
\item The reciprocal capacity $c^{-1}$ is the effective resistance of the rooted contact tree,
with leaf resistance $r_i^{-2}$.  Corona pruning is therefore exact series--parallel reduction.
\item Ordinary type-$A$ legs admit an exact path completion.  Their Schur complement replaces the
unit leaf resistance by $((l_i+1)r_i^2)^{-1}$ and preserves the one-defect law.
\item Quasi-Hamiltonian fusion and surface topology change only the coefficient of common rank,
producing the Euler term in \eqref{eq:intro-euler-capacity}.
\item The global dimension formula becomes an exact conservation law: one topological rank budget
minus positive leg fluctuations and positive fission-conormal losses.
\end{enumerate}

For honest graphs the theorem has a concise spectral consequence.  If $B$ is the adjacency matrix
and its eigenvalues are ordered decreasingly, then
\[
                         \lambda_2(B)\le2.
\]
If the spectral radius is below two the Cartan form is finite; if it equals two the eigenvalue two
is simple; if it exceeds two there is exactly one eigenvalue above two and none equal to two.  Thus
the known hyperbolic examples are the maximal indefinite possibility: higher-signature
Kac--Moody systems are excluded from type-$A$ genus-zero nonabelian Hodge diagrams.

The matrix-theory ancestors of the argument are the weighted-tree descriptions and inverse
$M$-matrix properties of ultrametric and generalized ultrametric matrices
\cite{MartinezMichonSanMartin1994,McDonaldEtAl1995,NabbenVarga1995}.  A recent treatment of shifted
generalized ultrametric matrices and tree-structured inverses is \cite{Nabben2026Acyclic}.  Those
works concern matrix recognition and inverses.  Exact conversion from Stokes contact to the Cartan
factorization \eqref{eq:intro-factor}, its capacity wall, and the full-leg passage form the local
input developed self-containedly below.  The new universal
contribution is the Euler--capacity identity \eqref{eq:euler-capacity-law}: it promotes scalar
rigidity to a polarized form, proves one-port inertia in every genus, and constrains every physical
rank reading.  Neither classical tree Gram decomposition, local one-defect law, nor known scalar
Euler-characteristic identity implies this global statement.  The sewing, variational, and
arbitrary-genus consequences are new; genus-zero conormal/Rees and spectral results are recovered
as collateral.  We state these distinctions because renaming an antecedent would not constitute a
new result.

The paper is organized as follows.  Section~\ref{sec:geometry} isolates the precise geometric
input.  Section~\ref{sec:factor} proves the one-defect factorization.  Section~\ref{sec:capacity}
proves the capacity trichotomy and root invariance.  Section~\ref{sec:corona} gives the electrical
corona recursion.  Section~\ref{sec:legs} treats arbitrary type-$A$ legs.  Sections
\ref{sec:conormal} and \ref{sec:dimension} give conormal/Rees and dimension formulations.
Section~\ref{sec:global} proves the arbitrary-genus Euler--capacity law by quasi-Hamiltonian
fusion, and
Section~\ref{sec:spectral} derives graph-theoretic consequences and examples.  The final section
records exact scope and falsifiers.

\section{Stokes contact and rooted fission data}\label{sec:geometry}

We separate the geometric input from the linear algebra.  This prevents a pairwise ultrametric
inequality from being silently upgraded to a simultaneous tree realization.

\subsection{Decorated cores}

Let $I$ be a nonempty finite set.  A \emph{decorated core} consists of a symmetric matrix
$B=(B_{ij})_{i,j\in I}$ and positive integers $r_i$.  In the Stokes application, $I$ indexes
Stokes circles, $r_i$ is the ramification order, and $B$ is the core matrix of the nonabelian
Hodge diagram.  Put
\begin{equation}\label{eq:Cartan-def}
                         C=2I-B.
\end{equation}
Diagonal entries of $B$ are retained; thus this notation includes generalized diagrams with
loops.  Define $K$ by \eqref{eq:intro-K}.  Both cases combine into
\begin{equation}\label{eq:B-contact}
            B_{ij}=\delta_{ij}+r_ir_j(K_{ij}-1).
\end{equation}

\begin{definition}\label{def:rooted-realization}
A \emph{rooted contact realization} of $K$ is a finite rooted tree $\Tcal$ whose leaves are
labelled by $I$, together with a real height $\kappa(v)$ at every retained vertex, strictly
increasing along each edge from a leaf toward the root, such that
\[
                         K_{ij}=\kappa(i\wedge j)
\]
for all $i,j\in I$.  Here $i\wedge i=i$, so leaf heights encode the diagonal contacts.
\end{definition}

For irregular classes, the realization is the common-truncation skeleton of the fission tree.
The edge and loop formulas for a Stokes circle give \eqref{eq:B-contact}; common Puiseux
truncations give the meet rule.  The following elementary reconstruction records why one common
tree, not a collection of pairwise witnesses, is present.

\begin{lemma}[Contact dendrogram]\label{lem:dendrogram}
Suppose a symmetric matrix $K$ satisfies $K_{ii}<K_{ij}$ for $i\ne j$ and, for every three
distinct indices, the two larger of $K_{ij},K_{ik},K_{jk}$ are equal.  Then $K$ has a rooted
contact realization, unique after unary vertices are suppressed.  Its clusters are the classes of
the relations
\[
 i\sim_a j\quad\Longleftrightarrow\quad i=j\ \text{or}\ K_{ij}\le a
\]
as $a$ ranges over the off-diagonal values of $K$.
\end{lemma}

\begin{proof}
Fix $a$.  Symmetry and reflexivity are clear.  If $i\sim_a j$ and $j\sim_a k$, then the stated
three-point condition forbids $K_{ik}$ from being strictly larger than both $K_{ij}$ and $K_{jk}$;
hence $K_{ik}\le a$.  Thus $\sim_a$ is an equivalence relation.  Its classes are nested as $a$
increases.  Make every non-singleton class a vertex, attach it to the least strictly larger class
that contains it, and attach the labelled singleton leaves at height $K_{ii}$.  Give a class born
at level $a$ height $a$.  The strict diagonal inequality makes pendant edges positive; deleting
repeated levels suppresses unary vertices.  The smallest cluster containing $i$ and $j$ is born at
$K_{ij}$, so its vertex is their meet.  Conversely, any rooted realization has exactly these
threshold clusters, proving uniqueness.
\end{proof}

\begin{proposition}[Geometric input]\label{prop:geometric-input}
The complementary-contact matrix of a finite type-$A$ irregular class has a rooted contact
realization.  For a genus-zero nonabelian Hodge diagram, any generic one-pole reading supplies such
a realization of its core, and the remaining vertices form ordinary type-$A$ legs.
\end{proposition}

\begin{proof}
At one irregular pole, the common truncations of the exponential factors form the fission tree.
Choose a level $\eta$ above all retained characteristic exponents.  On a common truncation branch,
let $R(u)>0$ be the cumulative ramification index at height $u$.  Telescoping the explicit loop and
edge formulae gives the complementary coarea coordinate
\begin{equation}\label{eq:contact-coarea}
 \kappa(v)=\eta-\int_{h(v)}^\eta\frac{du}{R(u)}.
\end{equation}
Profiles on two branches agree above their meet, so \eqref{eq:contact-coarea} is independent of
the leaf used to read a retained common-truncation vertex.  Formula
\eqref{eq:contact-coarea} is the summation-by-parts form of the rescaled source formula
\cite[Lemma~4.3]{Doucot2025Graphs}; the diagonal case is its loop analogue.  The diagonal formula reads
$K_{ii}=\kappa(i)$, and the pair formula reads $K_{ij}=\kappa(i\wedge j)$.  If $v$ is strictly below
$w$ on one branch, then
\[
 \kappa(w)-\kappa(v)=\int_{h(v)}^{h(w)}\frac{du}{R(u)}>0.
\]
Thus contact is strictly increasing toward the root.  This gives
Definition~\ref{def:rooted-realization}; equivalently one may apply Lemma~\ref{lem:dendrogram} to
the rescaled ultrametric relations proved from the explicit Stokes edge formulae in
\cite[Theorem~4.2]{Doucot2025Graphs}.

For a general algebraic connection on the affine line, the diagram is invariant under the
symplectic Weyl-algebra action \cite{Doucot2021Diagrams}.  The generic class reconstructed from the
enriched tree has all Stokes circles at infinity, while the enriched tree and diagram are invariant
under the symplectic action \cite[Theorems~1.4--1.6]{Doucot2026Basic}.  This supplies the required
one-pole reading.  The conjugacy-class data attach the ordinary type-$A$ legs used in the diagram
construction \cite{BoalchYamakawa2020,Doucot2026Basic}.  Only equality of the diagram is used here;
no claim that all weak representations have already been proved analytically isomorphic is needed.
\end{proof}

\begin{remark}
Proposition~\ref{prop:geometric-input} is the geometric boundary of the paper.  The remaining
proofs use only a rooted realization and the exact identity \eqref{eq:B-contact}.  Conversely, a
symmetric matrix satisfying an isolated collection of inequalities is not declared realizable by
an irregular connection.
\end{remark}

\subsection{Root extension and cut coordinates}

Choose $H>\max(1,\kappa(v))$ and attach a new root of height $H$.  Orient edges away from that
root.  For each edge $e$, let $A_e\subseteq I$ be its descendant leaves and let $\ell_e>0$ be the
drop in height along $e$.  Define
\begin{equation}\label{eq:cut-vectors}
 u_e=(r_i\one_{A_e}(i))_{i\in I},\qquad r=(r_i)_{i\in I},
 \qquad \lambda=H-1.
\end{equation}
The root extension is deliberately noncanonical.  Section~\ref{sec:capacity} will show that every
quantity claimed to be intrinsic is independent of it.

\section{The Stokes--Cartan one-defect factorization}\label{sec:factor}

\begin{lemma}[Root-path Gram identity]\label{lem:root-path}
For all leaves $i,j$,
\begin{equation}\label{eq:root-path}
 H-K_{ij}=\sum_e\ell_e\one_{A_e}(i)\one_{A_e}(j).
\end{equation}
Consequently the matrix $(H-K_{ij})$ is positive semidefinite.
\end{lemma}

\begin{proof}
An edge belongs to both root-to-leaf paths precisely when both leaves are descendants of that
edge.  The intersection of the two paths ends at $i\wedge j$.  Telescoping the positive height
drops along this common path gives $H-\kappa(i\wedge j)=H-K_{ij}$, proving
\eqref{eq:root-path}.  The right side is the Gram matrix
$\sum_e\ell_e a_ea_e^T$ with $a_e=(\one_{A_e}(i))_i$.
\end{proof}

\begin{theorem}[Local Stokes one-port factorization]\label{thm:one-defect}
For every rooted decorated Stokes core,
\begin{equation}\label{eq:one-defect-matrix}
 \boxed{\quad C=I+\sum_e\ell_eu_eu_e^T-\lambda rr^T.\quad}
\end{equation}
Equivalently, for every $x\in\RR^I$,
\begin{equation}\label{eq:one-defect-q}
 \boxed{\quad
 x^TCx=\sum_i x_i^2+
 \sum_e\ell_e\left(\sum_{i\in A_e}r_ix_i\right)^2
 -(H-1)\left(\sum_i r_ix_i\right)^2.\quad}
\end{equation}
\end{theorem}

\begin{proof}
By \eqref{eq:B-contact}, \eqref{eq:Cartan-def}, and Lemma~\ref{lem:root-path},
\begin{align*}
 C_{ij}
 &=\delta_{ij}+r_ir_j(1-K_{ij})\\
 &=\delta_{ij}+r_ir_j(H-K_{ij})-(H-1)r_ir_j\\
 &=\delta_{ij}+\sum_e\ell_e
   \bigl(r_i\one_{A_e}(i)\bigr)\bigl(r_j\one_{A_e}(j)\bigr)
   -\lambda r_ir_j.
\end{align*}
This is \eqref{eq:one-defect-matrix} entrywise.  Multiplying by $x_ix_j$ and summing gives
\eqref{eq:one-defect-q}.
\end{proof}

For later use, write
\begin{equation}\label{eq:mass-conormals}
 M(x)=r^Tx=\sum_i r_ix_i,
 \qquad \nu_e(x)=u_e^Tx=\sum_{i\in A_e}r_ix_i.
\end{equation}

\begin{corollary}[Canonical positive hyperplane]\label{cor:hyperplane}
The Cartan form is positive definite on
\[
                         \Hcal_r=\ker M.
\]
In particular $\indm(C)\le1$.
\end{corollary}

\begin{proof}
On $\Hcal_r$, the final term of \eqref{eq:one-defect-q} vanishes, while
$\sum_i x_i^2>0$ for $x\ne0$.  A quadratic form positive on a codimension-one subspace cannot have
a negative subspace of dimension two.
\end{proof}

The corollary is already the qualitative signature law.  The next section identifies the exact
wall and rules out a simultaneous zero and negative direction.

\section{Capacity, determinant, and the trichotomy}\label{sec:capacity}

Put
\begin{equation}\label{eq:P-c-Omega}
 P=I+\sum_e\ell_eu_eu_e^T,\qquad
 c=r^TP^{-1}r,\qquad
 \Omega=c^{-1}-\lambda.
\end{equation}
The matrix $P$ is positive definite, so $c>0$.

\begin{theorem}[Capacity normal form]\label{thm:capacity-normal}
For $x\in\RR^I$, set $m=M(x)$ and
\begin{equation}\label{eq:centered-z}
 z=x-\frac{m}{c}P^{-1}r.
\end{equation}
Then $r^Tz=0$, the decomposition is unique, and
\begin{equation}\label{eq:capacity-completion}
 \boxed{\quad x^TCx=z^TPz+m^2\Omega.\quad}
\end{equation}
\end{theorem}

\begin{proof}
Since $r^TP^{-1}r=c$, equation \eqref{eq:centered-z} gives $r^Tz=0$.  If
$x=z+aP^{-1}r$ with $r^Tz=0$, applying $r^T$ gives $a=m/c$, proving uniqueness.  Moreover
$z^TP(P^{-1}r)=z^Tr=0$.  Hence
\begin{align*}
 x^TCx
 &=x^TPx-\lambda(r^Tx)^2\\
 &=z^TPz+\frac{m^2}{c^2}r^TP^{-1}r-\lambda m^2\\
 &=z^TPz+m^2(c^{-1}-\lambda).
\end{align*}
\end{proof}

\begin{theorem}[Finite--affine--Lorentz trichotomy]\label{thm:trichotomy}
Exactly one of the following holds.
\begin{enumerate}[label=(\roman*),leftmargin=2.2em]
\item If $\Omega>0$, then $C$ is positive definite.
\item If $\Omega=0$, then $C$ is positive semidefinite,
\[
             \ker C=\spanr(P^{-1}r),\qquad \corank C=1.
\]
\item If $\Omega<0$, then $C$ is nonsingular and $\indm(C)=1$.
\end{enumerate}
In particular, $C$ has at most one nonpositive eigenvalue.  A negative and a zero eigenvalue can
never coexist.
\end{theorem}

\begin{proof}
In \eqref{eq:capacity-completion}, $z^TPz\ge0$, with equality exactly when $z=0$.
If $\Omega>0$, equality of the whole expression forces both $z=0$ and $m=0$, hence $x=0$.
If $\Omega=0$, the zero vectors are exactly those with $z=0$, namely the stated line.
If $\Omega<0$, the vector $P^{-1}r$ is negative because its centered part is zero and its mass is
$c>0$.  Corollary~\ref{cor:hyperplane} bounds the negative index by one, so it equals one.  A kernel
vector would have to be independent from this negative line or lie on it; the first possibility
would create a nonpositive plane and the second is excluded by strict negativity.  Equivalently,
the determinant formula below is nonzero.
\end{proof}

\begin{theorem}[Determinant wall]\label{thm:determinant}
One has
\begin{equation}\label{eq:determinant}
                   \boxed{\quad\det C=\det(P)c\Omega.\quad}
\end{equation}
Thus the affine wall is exactly $\Omega=0$, and it is crossed transversely in the scalar defect
$\lambda$.
\end{theorem}

\begin{proof}
The matrix determinant lemma gives
\[
 \det(P-\lambda rr^T)=\det(P)(1-\lambda r^TP^{-1}r)
 =\det(P)(1-\lambda c)=\det(P)c(c^{-1}-\lambda).
\]
\end{proof}

\begin{proposition}[Root-extension gauge invariance]\label{prop:root-invariance}
Extend the auxiliary root by length $t>0$.  Then
\begin{equation}\label{eq:root-update}
 P_t=P+t rr^T,\qquad \lambda_t=\lambda+t,\qquad
 c_t=\frac{c}{1+tc}.
\end{equation}
Consequently
\begin{equation}\label{eq:root-invariants}
 c_t^{-1}-\lambda_t=\Omega,\qquad
 \det(P_t)c_t=\det(P)c.
\end{equation}
In particular the trichotomy and determinant formula are intrinsic to the decorated diagram.
\end{proposition}

\begin{proof}
The new edge has descendant set $I$, so its cut vector is $r$, proving the first two updates.
Sherman--Morrison gives
\[
 P_t^{-1}=P^{-1}-\frac{tP^{-1}rr^TP^{-1}}{1+tc},
\]
and contraction with $r$ gives the capacity update.  Therefore
$c_t^{-1}=c^{-1}+t$, proving the first invariant.  The determinant lemma gives
$\det(P_t)=\det(P)(1+tc)$, which proves the second.
\end{proof}

\begin{remark}[Intrinsic formulation]
Choose any allowed root extension and compute $\Omega$.  Proposition~\ref{prop:root-invariance}
shows the result is independent of the choice.  Thus $\Omega$ is not a height normalization; it is
a scalar invariant of the Stokes--Cartan datum.  It is the signed distance between effective
resistance and global defect, as the next section makes literal.
\end{remark}

\section{Corona pruning and effective resistance}\label{sec:corona}

The term \emph{corona} will mean the terminal branches adjacent to a chosen internal vertex.  To
prune a corona is to eliminate all variables inside it while retaining the total current through
its stem.  The elimination is an exact Schur complement, not an asymptotic approximation.

Change variables from $x_i$ to leaf currents
\begin{equation}\label{eq:current-change}
                         m_i=r_ix_i.
\end{equation}
Then the positive core energy becomes
\begin{equation}\label{eq:electrical-energy}
 x^TPx=\sum_i\frac{m_i^2}{r_i^2}
       +\sum_e\ell_e\left(\sum_{i\in A_e}m_i\right)^2.
\end{equation}
Interpret $r_i^{-2}$ as a terminal resistance and $\ell_e$ as an edge resistance.  The current in
edge $e$ is the total leaf current below it.

\begin{lemma}[Parallel completion]\label{lem:parallel}
Let $R_1,\dots,R_s>0$ and $m_1+\cdots+m_s=M$.  Put
\[
 S=\sum_{a=1}^sR_a^{-1},\qquad R_{\parallel}=S^{-1}.
\]
Then
\begin{equation}\label{eq:parallel-completion}
 \sum_{a=1}^sR_am_a^2
 =\sum_{a=1}^sR_a\left(m_a-\frac{R_a^{-1}}{S}M\right)^2
  +R_{\parallel}M^2.
\end{equation}
\end{lemma}

\begin{proof}
Expand the squares.  The cross terms sum to $-2M^2/S$ and the constant terms to $M^2/S$.
Rearrangement gives \eqref{eq:parallel-completion}.  Positivity also shows that the unique
energy-minimizing split is proportional to conductance.
\end{proof}

\begin{definition}[Recursive resistance]\label{def:recursive-R}
Assign a leaf $i$ the resistance $R_i=r_i^{-2}$.  If an internal vertex $v$ has children $w$, and
the edge from $v$ to $w$ has length $\ell_{vw}$, define
\begin{equation}\label{eq:resistance-recursion}
                  R_v=\left(\sum_{w\prec v}
                    \frac{1}{\ell_{vw}+R_w}\right)^{-1}.
\end{equation}
The value at the chosen root is $R_{\mathrm{eff}}$.
\end{definition}

\begin{theorem}[Corona-capacity law]\label{thm:corona-capacity}
Successive pruning by \eqref{eq:resistance-recursion} gives
\begin{equation}\label{eq:min-energy}
 \min_{\sum_i m_i=M}x^TPx=R_{\mathrm{eff}}M^2,
 \qquad R_{\mathrm{eff}}=c^{-1}.
\end{equation}
Moreover
\begin{equation}\label{eq:Omega-resistance}
                         \Omega=R_{\mathrm{eff}}-\lambda.
\end{equation}
A root extension adds the same series resistance $t$ to both terms in
\eqref{eq:Omega-resistance}.  Thus root invariance is the series-cancellation law
\[
 (R_{\mathrm{eff}}+t)-(\lambda+t)=R_{\mathrm{eff}}-\lambda.
\]
\end{theorem}

\begin{proof}
At a terminal corona, the resistance from a child subtree is its effective resistance plus the
stem resistance $\ell_{vw}$ in series.  Lemma~\ref{lem:parallel} replaces all child branches by
their parallel resistance and a sum of positive fluctuation squares.  Induction toward the root
proves the first minimization formula.

On the other hand, Theorem~\ref{thm:capacity-normal} applied to the positive form alone gives
\[
 x^TPx=z^TPz+M^2/c
\]
under the constraint $r^Tx=M$.  Its unique minimum is $M^2/c$.  Comparing the two minima gives
$R_{\mathrm{eff}}=c^{-1}$, and \eqref{eq:Omega-resistance} follows from definition.
The final assertion is ordinary series addition together with
Proposition~\ref{prop:root-invariance}.
\end{proof}

\begin{corollary}[Local-to-global computation]\label{cor:message-passing}
The signature phase of a rooted core is determined by one bottom-up pass through the fission tree:
compute $R_v$ by \eqref{eq:resistance-recursion} and compare $R_{\mathrm{eff}}$ with $H-1$.
Every pruned corona contributes only positive fluctuation squares; no local pruning step can create
a second negative direction.
\end{corollary}

\begin{remark}[Vanishing-cycle language]
Formula \eqref{eq:parallel-completion} is an algebraic elimination law for a terminal cluster.  It
is reasonable to view the discarded positive fluctuation coordinates as quadratic
``vanishing-cycle channels.''  We do not identify them with sheaf-theoretic vanishing cycles, nor
do we claim that $\ell_e$ is a characteristic-cycle multiplicity.  Such an identification would
require an additional microlocal comparison theorem not proved here.
\end{remark}

\section{Type-A legs preserve the one-defect law}\label{sec:legs}

Attach to core vertex $i$ an ordinary path leg of length $l_i\ge0$.  Write its variables as
$y_{i,1},\dots,y_{i,l_i}$, and set
\[
 y_{i,0}=x_i,\qquad y_{i,l_i+1}=0,\qquad
 d_{i,k}=y_{i,k}-y_{i,k+1}\quad(0\le k\le l_i).
\]
Then $\sum_{k=0}^{l_i}d_{i,k}=x_i$.

\begin{lemma}[Path completion]\label{lem:path-completion}
For $d_0+\cdots+d_l=x$,
\begin{equation}\label{eq:path-completion}
 \sum_{k=0}^l d_k^2
 =\sum_{k=0}^l\left(d_k-\frac{x}{l+1}\right)^2+\frac{x^2}{l+1}.
\end{equation}
\end{lemma}

\begin{proof}
Expand the right side and use $\sum d_k=x$.
\end{proof}

Let $\widehat C$ be the full Cartan matrix.  The ordinary type-$A$ edge coupling gives
\begin{equation}\label{eq:full-before-completion}
 \widehat q(x,y)=
 \sum_i\sum_{k=0}^{l_i}d_{i,k}^2
 +\sum_e\ell_e\nu_e(x)^2-\lambda M(x)^2.
\end{equation}
Indeed, the path energy $\sum_kd_{i,k}^2$ equals the leg Cartan energy together with the core
square $x_i^2$ that occurs in \eqref{eq:one-defect-q}.

\begin{theorem}[Full-leg conservation]\label{thm:full-leg}
For arbitrary finite type-$A$ legs,
\begin{align}\label{eq:full-completed}
 \widehat q(x,y)
 &=\sum_i\sum_{k=0}^{l_i}
     \left(d_{i,k}-\frac{x_i}{l_i+1}\right)^2
   +\sum_i\frac{x_i^2}{l_i+1}\notag\\
 &\quad+\sum_e\ell_e\nu_e(x)^2-\lambda M(x)^2.
\end{align}
Consequently $\widehat C$ is positive definite, positive semidefinite of corank one, or has
exactly one negative eigenvalue.  In particular
\begin{equation}\label{eq:full-index}
                 \indm(\widehat C)\le1,\qquad
                 \dim\ker\widehat C\le1.
\end{equation}
\end{theorem}

\begin{proof}
Apply Lemma~\ref{lem:path-completion} to every leg in
\eqref{eq:full-before-completion}.  The fluctuation squares form a positive definite complement
to the variables $x$.  The remaining Schur form is
\begin{equation}\label{eq:leg-Schur}
 S=D+\sum_e\ell_eu_eu_e^T-\lambda rr^T,\qquad
 D=\diag\left(\frac1{l_i+1}\right)>0.
\end{equation}
Replace $I$ by $D$ in Sections~\ref{sec:factor}--\ref{sec:capacity}.  The positive matrix
$P_L=D+\sum_e\ell_eu_eu_e^T$ is still positive definite, so the same centered decomposition and
trichotomy apply.  Completing the legs is a congruence by a positive block and therefore preserves
the nonpositive inertia of $S$.
\end{proof}

\begin{corollary}[Leg capacity]\label{cor:leg-capacity}
For the full diagram, use terminal resistances
\begin{equation}\label{eq:leg-resistance}
                         R_i=\frac{1}{(l_i+1)r_i^2}.
\end{equation}
The recursion \eqref{eq:resistance-recursion} computes the full capacity wall.  In particular,
lengthening a leg changes only a positive terminal resistance and never increases the maximal
negative index.
\end{corollary}

\begin{proof}
Under $m_i=r_ix_i$, the diagonal term in \eqref{eq:leg-Schur} is
$m_i^2/((l_i+1)r_i^2)$.  Apply Theorem~\ref{thm:corona-capacity}.
\end{proof}

\section{Principal conormals and Rees stability}\label{sec:conormal}

The cut form has a direct microlocal interpretation that is exact at the level of linear
conormals.  For an edge $e$, define the cluster hyperplane
\begin{equation}\label{eq:cluster-hyperplane}
             \Hcal_e=\{x\in\RR^I:\nu_e(x)=0\}.
\end{equation}
Its conormal line is generated by the constant covector
\begin{equation}\label{eq:principal-conormal}
                         d\nu_e=u_e.
\end{equation}
The family $\{A_e\}$ is laminar: two descendant sets are disjoint or one contains the other.

\begin{theorem}[Principal-conormal Hessian law]\label{thm:conormal-Hessian}
Let $q_+(x)=x^TPx$.  Its Hessian satisfies
\begin{equation}\label{eq:Hessian-law}
 \frac12\operatorname{Hess}(q_+)
 =I+\sum_e\ell_e\,d\nu_e\otimes d\nu_e.
\end{equation}
For the full diagram, $I$ is replaced by
$D=\diag((l_i+1)^{-1})$.  Thus the complete local Stokes contribution is the weighted Gram norm of
the principal conormals of the laminar fission cuts.  The sole negative Hessian channel is
$-\lambda\,dM\otimes dM$.
\end{theorem}

\begin{proof}
This is the polarization of \eqref{eq:one-defect-q} and
\eqref{eq:full-completed}; $d\nu_e=u_e$ and $dM=r$.
\end{proof}

We next record the exact Rees-type stability that follows from this decomposition.  Assign an
indeterminate $s_e$ to each fission edge and define the multiscale family
\begin{equation}\label{eq:Rees-family}
 q_{\boldsymbol s}(x)=x^TDx+
   \sum_e s_e\ell_e\nu_e(x)^2-\lambda M(x)^2.
\end{equation}
At $s_e=1$ this is the reduced full Cartan form.  Setting selected $s_e$ to zero contracts their
graded conormal channels; scaling $s_e=t^{w_e}$ records a weighted contact filtration.  We call
\eqref{eq:Rees-family} the \emph{quadratic multi-Rees family} of the fission filtration.  This
terminology refers only to the explicit filtered quadratic family, not to a claimed equivalence of
categories of Stokes-filtered sheaves.

\begin{corollary}[Rees one-defect stability]\label{cor:Rees-stability}
For every specialization $s_e\ge0$, the form $q_{\boldsymbol s}$ has negative index at most one and
nullity at most one.  It is positive on $\ker M$.  Its phase is determined by
\begin{equation}\label{eq:Rees-gap}
 \Omega(\boldsymbol s)=
 \left[r^T\left(D+\sum_e s_e\ell_eu_eu_e^T\right)^{-1}r\right]^{-1}
 -\lambda.
\end{equation}
\end{corollary}

\begin{proof}
The specialized positive part remains positive definite because $D>0$ and $s_e\ell_e\ge0$.
The proofs of Corollary~\ref{cor:hyperplane} and Theorem~\ref{thm:trichotomy} apply verbatim.
\end{proof}

\begin{remark}
Equation \eqref{eq:Hessian-law} is a principal-conormal statement in a finite linear model.  It
does not identify $\ell_e$ with a multiplicity in the characteristic cycle of an irregular
holonomic $\mathcal D$-module.  Establishing such a statement would require compatibility between
the irregular Riemann--Hilbert correspondence, microlocal specialization, and the fission-contact
filtration.  The present theorem supplies a rigid numerical target for that future comparison.
\end{remark}

\section{The wild dimension conservation law}\label{sec:dimension}

Let $\widehat d=(x,y)$ be a dimension vector for the full diagram, written in the coordinates of
Section~\ref{sec:legs}.  The diagram dimension formula is
\begin{equation}\label{eq:dimension-known}
                \dim\Mcal_B=2-\widehat d^T\widehat C\widehat d
\end{equation}
for the corresponding wild character variety, under the usual nonempty/stability hypotheses
\cite{Doucot2021Diagrams}.

\begin{theorem}[Rank--conormal dimension conservation]\label{thm:dimension-law}
With $d_{i,k}=y_{i,k}-y_{i,k+1}$ as above,
\begin{align}\label{eq:dimension-law}
 \dim\Mcal_B
 &=2+\lambda M(x)^2
 -\sum_i\sum_{k=0}^{l_i}
     \left(d_{i,k}-\frac{x_i}{l_i+1}\right)^2\notag\\
 &\quad-\sum_i\frac{x_i^2}{l_i+1}
 -\sum_e\ell_e\nu_e(x)^2.
\end{align}
This equality is independent of the auxiliary root extension.  It expresses dimension as one
global ramified-rank square budget minus positive leg fluctuations, positive terminal costs, and
positive principal-conormal defects.
\end{theorem}

\begin{proof}
Substitute \eqref{eq:full-completed} into \eqref{eq:dimension-known}.  If the root is extended by
$t$, the new root edge has $A_e=I$ and hence contributes $tM(x)^2$ to the final positive sum, while
$\lambda M(x)^2$ increases by the same amount.  The two changes cancel.
\end{proof}

\begin{corollary}[Relaxed dimension obstruction]\label{cor:dimension-obstruction}
For fixed total ramified rank $M$ and fixed root gauge, the real relaxation of the dimension
functional is maximized by the unique conductance-proportional current split at every corona and by
constant increments along every type-$A$ leg.  Every departure from these balances has an explicit
positive square cost.  An integral dimension vector attains this relaxed bound only when the
balanced profile is integral.
\end{corollary}

\begin{proof}
Apply Lemmas~\ref{lem:parallel} and \ref{lem:path-completion} to
\eqref{eq:dimension-law}.  Equality in either completion gives the stated unique balance.
\end{proof}

This corollary is a useful distinction between the universal law and a sign estimate.  The theorem
does not merely say that certain terms are nonnegative; it locates every loss channel and gives its
equality condition.

\section{The Euler--capacity law on a punctured curve}\label{sec:global}

We now pass from one-puncture atoms to a connected curve of arbitrary genus.  This step uses only
the quasi-Hamiltonian dimension formula, but it changes the scope of the result: the genus-zero
Cartan theorem becomes one boundary case of a topological index law.

Let $X$ be a compact connected curve of genus $g$, let
$D=\{p_1,\ldots,p_m\}$ with $m\ge1$, and put $G=\mathrm{GL}_M(\CC)$.  At $p_a$ choose an irregular
type $\mathcal Q_a$, with formal centralizer $H_a$, and a conjugacy class
$\mathcal C_a\subset H_a$, and let $\mathcal A_a$ be its local fission/Stokes space.  Fixing the
combinatorial marking makes the following dimensions quadratic functions of the marked
multiplicities:
\[
 h_a=\dim H_a,\qquad k_a=\dim\mathcal C_a,\qquad
 s_a=\dim\mathcal A_a-M^2-h_a.
\]
For an unramified irregular type in a Cartan subalgebra,
$s_a=\sum_{\alpha\in\mathcal R(G)}\deg(\alpha\circ\mathcal Q_a)$; the intrinsic definition above
also covers the ramified local Stokes spaces used here.  Thus the local space has dimension
$M^2+h_a+s_a$ \cite{Boalch2014Geometry,BoalchYamakawa2015Twisted}.  Define the global numerical index by
\begin{equation}\label{eq:global-index-definition}
 Q_{X,D}:=2-\dim\Mcal_{X,D}
\end{equation}
whenever the stable reduction is nonempty, and otherwise by the same quadratic dimension
polynomial.  Define the one-puncture atom
\begin{equation}\label{eq:local-index-atom}
                         Q_a:=M^2+h_a-s_a-k_a.
\end{equation}
The marked type-$A$ diagram construction identifies \eqref{eq:local-index-atom} with the full local
Cartan quadratic form, including its conjugacy legs
\cite{BoalchYamakawa2020,Doucot2021Diagrams}.

\begin{proposition}[Exact fusion identity]\label{prop:fusion-identity}
On the rank fibre product \eqref{eq:intro-fibre-product},
\begin{equation}\label{eq:fusion-identity}
             \boxed{Q_{X,D}=\sum_{a=1}^mQ_a+2(1-g-m)M^2.}
\end{equation}
\end{proposition}

\begin{proof}
The quasi-Hamiltonian Stokes representation space has dimension
\[
 \dim\operatorname{Hom}_{\mathbb S}
 =(2g-2)M^2+\sum_a(M^2+h_a+s_a).
\]
Reduction at $\prod_a\mathcal C_a$ subtracts
$2(\sum_ah_a-\dim Z(G))$ and adds $\sum_ak_a$.  Since
$\dim Z(\mathrm{GL}_M)=1$,
\begin{equation}\label{eq:global-dimension-expanded}
 \dim\Mcal_{X,D}
 =2+(2g-2+m)M^2+\sum_a(s_a+k_a-h_a).
\end{equation}
Therefore
\[
 Q_{X,D}=(2-2g-m)M^2+\sum_a(h_a-s_a-k_a).
\]
Subtracting the sum of \eqref{eq:local-index-atom} gives
$2(1-g-m)M^2$, proving \eqref{eq:fusion-identity}.  Because every term is a quadratic polynomial
in the marking, equality on integral multiplicities polarizes to the real marked lattice.
\end{proof}

For each $a$, apply the local full-leg theorem.  Write $R_a:V_a\to\RR$ for total ramified rank.
There is a positive-definite form $P_a$, a capacity vector $v_a$ normalized by $R_a(v_a)=1$, and
an intrinsic scalar $\Omega_a$ such that every $x_a\in V_a$ decomposes uniquely as
\begin{equation}\label{eq:local-capacity-splitting}
 x_a=z_a+R_a(x_a)v_a,qquad R_a(z_a)=0,qquad
 Q_a(x_a)=P_a(z_a)+\Omega_aR_a(x_a)^2.
\end{equation}
Concretely, $P_a$ is the positive Schur form in Theorem~\ref{thm:full-leg},
$v_a=P_a^{-1}R_a^*/c_a$, $c_a=R_aP_a^{-1}R_a^*$, and
$\Omega_a=c_a^{-1}-\lambda_a$.  Put
\begin{equation}\label{eq:local-capacity-defect}
                         \delta_a:=2-\Omega_a.
\end{equation}
This normalization is forced by gluing: a one-puncture sphere has Euler characteristic two, so its
global coefficient is $2-\delta_a=\Omega_a$.

\begin{theorem}[Euler--capacity and one-port index law]\label{thm:stokes-hodge}
For every $x=(x_a)\in V_{X,D}$ of common rank $M$, one has the orthogonal conservation law
\begin{equation}\label{eq:euler-capacity-law}
 \boxed{\quad Q_{X,D}(x)=\sum_{a=1}^mP_a(z_a)+\Omega_{X,D}M^2,\qquad
 \Omega_{X,D}=\chi(X)-\sum_{a=1}^m\delta_a
 =\chi(X)+\sum_{a=1}^m(\Omega_a-2).\quad}
\end{equation}
Thus the sole scalar channel obeys a Gauss--Bonnet form: global Euler characteristic minus the sum
of renormalized local capacity defects.
Consequently:
\begin{enumerate}[label=(\roman*),leftmargin=2.2em]
\item $Q_{X,D}$ is positive definite on the physical rank-zero hyperplane;
\item $\indm(Q_{X,D})+\dim\ker(Q_{X,D})\le1$;
\item $Q_{X,D}$ is positive definite, positive semidefinite with a one-dimensional kernel, or
Lorentzian according as $\Omega_{X,D}>0$, $=0$, or $<0$;
\item in the affine case the null line is spanned by the simultaneous capacity ray
$(v_1,\ldots,v_m)$.
\end{enumerate}
\end{theorem}

\begin{proof}
Insert \eqref{eq:local-capacity-splitting} into Proposition~\ref{prop:fusion-identity}.  The common
rank condition gives
\[
 Q_{X,D}(x)=\sum_aP_a(z_a)
 +\left(\sum_a\Omega_a+2(1-g-m)\right)M^2.
\]
Since $\chi(X)=2-2g$, the coefficient is
$\chi(X)+\sum_a(\Omega_a-2)$.  The first summand is positive definite in all centered variables.
Thus the displayed splitting is an orthogonal sum of a positive-definite space and one scalar
line, which proves all inertia statements and identifies the null line.
\end{proof}

\begin{corollary}[Variational characterization]\label{cor:variational-gap}
The Euler--capacity coefficient is intrinsic and is the constrained ground-state energy
\begin{equation}\label{eq:global-variational-gap}
 \boxed{\quad \Omega_{X,D}=min\{Q_{X,D}(x):x\in V_{X,D},\ M(x)=1\}.\quad}
\end{equation}
The minimum is attained uniquely at the simultaneous capacity ray $(v_1,\ldots,v_m)$.  More
generally, the minimum at common rank $M$ is $\Omega_{X,D}M^2$.  Locally,
\begin{equation}\label{eq:local-defect-variational}
 \Omega_a=\min_{R_a(x)=1}Q_a(x),\qquad
 \delta_a=2-\min_{R_a(x)=1}Q_a(x).
\end{equation}
\end{corollary}

\begin{proof}
In \eqref{eq:euler-capacity-law}, the positive sum is minimized uniquely by $z_a=0$ for every $a$.
The stated values and minimizers follow.  The local assertion is the case $g=0,m=1$.
\end{proof}

\begin{corollary}[Physical-rank covectors]\label{cor:physical-rank}
Let a fixed genus-zero nonabelian Hodge quadratic form $q$ be presented by any admissible
type-$A$ Stokes reading $\alpha$, and let $R_\alpha$ be rank in that reading.  Then
\begin{equation}\label{eq:rank-kernel-positive}
                  q(x)>0\qquad(0\ne x\in\ker R_\alpha).
\end{equation}
If $q$ is nonsingular Lorentzian, this is equivalently
\begin{equation}\label{eq:dual-timelike-rank}
                         R_\alpha q^{-1}R_\alpha^T<0.
\end{equation}
Thus Fourier--Laplace or another basic diagram-preserving transformation may change numerical
rank, but every geometric rank covector remains dual-timelike for the invariant form.  No such
statement holds for arbitrary positive covectors.
\end{corollary}

\begin{proof}
Apply Theorem~\ref{thm:stokes-hodge} to the local atoms in reading $\alpha$; common rank zero removes
the sole scalar term in \eqref{eq:euler-capacity-law}, leaving a nonzero sum of positive-definite
forms.  In the Lorentz case, the equivalence with \eqref{eq:dual-timelike-rank} is the standard
one-codimensional Schur-complement criterion, obtained by centering along $q^{-1}R_\alpha^T$.
\end{proof}

\begin{corollary}[Fusion and handle laws]\label{cor:fusion-handles}
If the positive local atoms are first fused before the global rank defect is inserted, their
effective capacities satisfy
\begin{equation}\label{eq:fused-capacity}
                         c_{\mathrm{fus}}^{-1}=\sum_{a=1}^m c_a^{-1}.
\end{equation}
Increasing the genus by one changes no centered term and sends
$\Omega_{X,D}$ to $\Omega_{X,D}-2$.
\end{corollary}

\begin{proof}
At common rank $M$, minimization of $\sum_aP_a$ gives
$M^2\sum_ac_a^{-1}$, proving \eqref{eq:fused-capacity}.  The handle statement follows immediately
from $\chi(X)=2-2g$.
\end{proof}

\begin{corollary}[Topological monotonicity and sewing]\label{cor:topological-monotonicity}
For fixed local data, increasing genus moves the phase monotonically through the ordered sequence
finite, affine, Lorentz, and the affine phase can occur at at most one integer genus.  If two
punctured curves are joined by connected sum away from their marked points, then
\begin{equation}\label{eq:connected-sum-gap}
 \Omega_{X\mathbin\#Y,D\sqcup E}=\Omega_{X,D}+\Omega_{Y,E}-2.
\end{equation}
\end{corollary}

\begin{proof}
The genus step is $\Omega\mapsto\Omega-2$ by Corollary~\ref{cor:fusion-handles}.  For connected sum,
use $\chi(X\#Y)=\chi(X)+\chi(Y)-2$ and add the two disjoint sets of local capacity defects in
Theorem~\ref{thm:stokes-hodge}.
\end{proof}

\begin{corollary}[Genus-zero diagram law]\label{thm:global}
Let $\widehat C$ be the symmetric Cartan matrix of a type-$A$ genus-zero nonabelian Hodge diagram.
Then
\begin{equation}\label{eq:global-index}
 \indm(\widehat C)+\dim\ker\widehat C\le1.
\end{equation}
Its phase, affine null line, determinant wall, and dimension may be computed in any generic
one-puncture reading.
\end{corollary}

\begin{proof}
Diagram invariance permits a generic symplectic reading of the same core at infinity
\cite{Doucot2021Diagrams,Doucot2026Basic}.  Proposition~\ref{prop:geometric-input} supplies its
rooted fission-contact realization, and Theorem~\ref{thm:full-leg} gives
\eqref{eq:local-capacity-splitting}.  Apply Theorem~\ref{thm:stokes-hodge} with $g=0,m=1$.
\end{proof}

\begin{corollary}[Fourier invariance of the wall]\label{cor:Fourier-wall}
The sign of the one-puncture capacity gap and the finite/affine/Lorentz phase are invariant under
Fourier--Laplace and symplectic Weyl-algebra transformations preserving the genus-zero diagram.
\end{corollary}

\begin{proof}
The sign is the inertia of the invariant Cartan matrix by Theorem~\ref{thm:stokes-hodge}.
\end{proof}

\begin{remark}[Irregular Riemann--Hilbert and spectral-curve scope]
The irregular Riemann--Hilbert correspondence transports the same numerical index to generalized
monodromy and Stokes data.  Under Hiroe's additional hypotheses, its value at an integral
multiplicity vector also equals the Euler characteristic of the normalized irregular spectral
curve \cite{Hiroe2019Rigidity}.  Equation~\eqref{eq:euler-capacity-law} is a quadratic refinement:
it determines the whole real polarization and its unique possible negative channel.  It is not a
new proof of irregular Riemann--Hilbert, and it neither reconstructs a Stokes local system from
$\Omega_{X,D}$ nor identifies the positive squares with sheaf-theoretic vanishing cycles.
\end{remark}

\section{Graph spectra and Kac--Moody exclusion}\label{sec:spectral}

Suppose the full diagram is an honest graph: $B$ is a real symmetric adjacency matrix with zero
diagonal and nonnegative off-diagonal entries.  Let
$\lambda_1(B)\ge\cdots\ge\lambda_n(B)$ be its eigenvalues.

\begin{theorem}[Second-eigenvalue obstruction]\label{thm:lambda2}
If $n\ge2$ and $B$ is a type-$A$ genus-zero nonabelian Hodge graph, then
\begin{equation}\label{eq:lambda2}
                         \boxed{\lambda_2(B)\le2.}
\end{equation}
More precisely:
\begin{enumerate}[label=(\roman*),leftmargin=2.2em]
\item if $\lambda_1(B)<2$, then $2I-B$ is positive definite;
\item if $\lambda_1(B)=2$, then $2$ is a simple eigenvalue and $2I-B$ is positive semidefinite;
\item if $\lambda_1(B)>2$, then exactly one eigenvalue of $B$ exceeds $2$ and $2$ is not an
eigenvalue.
\end{enumerate}
\end{theorem}

\begin{proof}
The eigenvalues of $2I-B$ are $2-\lambda_i(B)$.  Apply
Theorem~\ref{thm:global}: its negative index is the number of adjacency eigenvalues above two, and
its nullity is the multiplicity of the eigenvalue two.  Formula \eqref{eq:global-index} gives all
claims.
\end{proof}

\begin{corollary}[Kac--Moody signature bound]\label{cor:Kac}
No symmetric generalized Cartan matrix of negative index at least two can be realized by a
type-$A$ genus-zero nonabelian Hodge graph.  Hyperbolic/Lorentzian type is the maximal indefinite
type in this class.
\end{corollary}

The statement is deliberately restricted to the class in Theorem~\ref{thm:global}.  It does not
exclude higher-signature generalized Cartan matrices from other quiver varieties, other genera,
or constructions not admitting the rooted type-$A$ fission reading.

\begin{example}[A sharp twisted graph]\label{ex:twisted-triangle}
Dou\c cot constructs a nonabelian Hodge graph from a twisted irregular class whose three edge
multiplicities are $6,3,4$ \cite[Proposition~3.3]{Doucot2025Graphs}.  In a compatible ordering,
\[
 B=\begin{pmatrix}0&6&4\\6&0&3\\4&3&0\end{pmatrix},\qquad
 C=\begin{pmatrix}2&-6&-4\\-6&2&-3\\-4&-3&2\end{pmatrix}.
\]
For $\one=(1,1,1)^T$,
\[
                         \one^TC\one=6-2(6+4+3)=-20.
\]
Also $\det C=-258$ and $\operatorname{tr}C=6$, independently confirming that the inertia is
$(2,1,0)$.
Thus the universal theorem cannot be strengthened to positivity.  Its conclusion is sharp:
$C$ has exactly one negative eigenvalue and no kernel.  This graph is also useful because it is a
nonabelian Hodge graph that is not an untwisted fission graph, so the phenomenon is not confined to
complete multipartite cores.
\end{example}

\begin{example}[Affine wall]
Whenever a diagram lies on $R_{\mathrm{eff}}=H-1$, its null vector is explicitly
$P^{-1}r$.  In electrical variables, this is the unique conductance-balanced unit-current profile.
The simplicity of the affine eigenvalue is therefore not an appeal to a Dynkin classification; it
is forced directly by strict convexity of the positive Stokes energy.
\end{example}

\begin{example}[Why a second defect is impossible]
Assume two independent vectors had nonpositive Cartan energy.  A nontrivial linear combination of
them has zero total ramified rank.  On that hyperplane the energy is the strictly positive sum of
terminal and cut squares in \eqref{eq:one-defect-q}, a contradiction.  This two-line argument is
the geometric core of the exclusion, while capacity supplies its sharp boundary behavior.
\end{example}

\section{Sharpness, falsifiers, and limits}\label{sec:scope}

We finish by listing what would disprove or extend the theorem.  This is part of the mathematical
statement: a universal-law claim without explicit falsifiers is not auditable.

\subsection{Statements ruled out}

The stronger conjecture ``all nonabelian Hodge Cartan forms are positive semidefinite'' is false by
Example~\ref{ex:twisted-triangle}.  The correct boundary is negative index one.  The theorem also
rules out a weaker but previously plausible phenomenon: a Lorentz direction together with an
independent affine direction.  Capacity completion shows that this cannot occur.

The statement ``every positive covector is a physical rank port'' is also false.  For the twisted
triangle of Example~\ref{ex:twisted-triangle}, direct inversion gives
\[
 \begin{aligned}
 e_1^TC^{-1}e_1&=\frac5{258}>0,\\
 (3,2,3)C^{-1}(3,2,3)^T&=-\frac{245}{86}<0,\\
 (2,1,4)C^{-1}(2,1,4)^T&=-\frac{104}{129}<0.
 \end{aligned}
\]
The latter two covectors are physical ranks in two symplectic readings; $e_1$ is not.  Thus the
rank fibre product in \eqref{eq:intro-fibre-product} is geometric content, not replaceable by an
arbitrary choice of a positive linear functional.

\subsection{Direct falsifiers}

Any one of the following would invalidate the stated scope.
\begin{enumerate}[label=(\arabic*),leftmargin=2.2em]
\item A marked type-$A$ local Stokes atom whose complementary contacts do not admit the asserted
simultaneous rooted realization, or whose completed local index has two negative directions.
\item A quasi-Hamiltonian wild character variety for $\mathrm{GL}_M$ whose dimension violates
\eqref{eq:global-dimension-expanded}; equivalently, a handle or fusion operation that changes a
centered local mode rather than only common rank.
\item A nonzero compatible deformation of common rank zero with
$Q_{X,D}\le0$.
\item A failure of diagram preservation under the generic symplectic reading used in
Corollary~\ref{thm:global}.
\end{enumerate}
The first and fourth items are the local geometric inputs; the second is the published
quasi-Hamiltonian dimension formula; the third is a direct consequence and a concrete falsification
test.  Conversely, the theorem is not an analytic existence criterion: a quadratic form satisfying
the one-defect law need not come from an irregular connection.

\subsection{Boundaries of interpretation}

The conormals in Theorem~\ref{thm:conormal-Hessian} are conormals of explicit cluster hyperplanes in
dimension-vector space.  The corona reduction is an exact Schur complement.  The Rees family is a
filtered deformation of the quadratic form.  None of these statements, alone, computes
microsupports, Stokes sheaf categories, characteristic-cycle coefficients, or vanishing-cycle
multiplicities.  Those are natural next questions because the present law predicts that any
compatible categorification must preserve a single global negative channel while decomposing all
local fission channels positively.

\subsection{Novelty boundary}

Weighted-tree Gram formulae, ultrametric matrices, inverse $M$-matrices, the Stokes diagram, and its
dimension formula are prior results.  So are quasi-Hamiltonian fusion and, under additional
hypotheses, equality of the scalar rigidity index with a spectral-curve Euler characteristic.
The claim here is the previously unstated quadratic refinement: the Stokes normalization
\eqref{eq:intro-K} makes every local centered direction positive, while the exact fusion identity
shows that Euler topology and all local capacities meet in one common-rank coefficient
\eqref{eq:euler-capacity-law}.  This determines full inertia, not merely the value at one dimension
vector.  Corona pruning, type-$A$ legs, Rees scaling, Fourier readings, and the Kac--Moody spectrum
are controlled shadows of that law.  Removing the rooted Stokes realization or the physical-rank
compatibility condition destroys the conclusion.

\section{Conclusion}

The universal law is \eqref{eq:euler-capacity-law}.  A connected irregular curve has many local
Stokes channels but only one common physical-rank channel.  Fission contacts, type-$A$ markings,
corona elimination, and Rees deformation contribute positive centered energy.  Fusion and handles
cannot create a second defect; they alter only the coefficient of rank, by the Euler correction.
That is why the full index is at most Lorentzian in every genus.

The genus-zero Stokes--Cartan factorization, determinant wall, effective-resistance recursion,
Fourier wall, dimension conservation law, and adjacency bound $\lambda_2\le2$ are therefore
collateral manifestations of one Euler--capacity index law.  The remaining hard problem is
categorical: construct an irregular Riemann--Hilbert/Rees comparison whose vanishing-cycle or
characteristic-cycle pairing realizes the positive principal-conormal summands.  The theorem here
gives the exact numerical polarization, sign constraint, and fusion law that such a comparison must
recover.

\begin{thebibliography}{99}

\bibitem{BiquardBoalch2004}
O.~Biquard and P.~Boalch,
\emph{Wild non-abelian Hodge theory on curves},
Compos. Math. \textbf{140} (2004), no.~1, 179--204.
\href{https://doi.org/10.1112/S0010437X03000010}{doi:10.1112/S0010437X03000010}.

\bibitem{Boalch2008Kac}
P.~Boalch,
\emph{Irregular connections and Kac--Moody root systems},
arXiv:0806.1050 (2008).
\href{https://arxiv.org/abs/0806.1050}{arXiv:0806.1050}.

\bibitem{Boalch2012Simply}
P.~Boalch,
\emph{Simply-laced isomonodromy systems},
Publ. Math. Inst. Hautes \'{E}tudes Sci. \textbf{116} (2012), 1--68.
\href{https://doi.org/10.1007/s10240-012-0040-2}{doi:10.1007/s10240-012-0040-2}.

\bibitem{Boalch2014Geometry}
P.~Boalch,
\emph{Geometry and braiding of Stokes data; fission and wild character varieties},
Ann. of Math. (2) \textbf{179} (2014), no.~1, 301--365.
\href{https://doi.org/10.4007/annals.2014.179.1.5}{doi:10.4007/annals.2014.179.1.5}.

\bibitem{Boalch2025Counting}
P.~Boalch,
\emph{Counting the fission trees and nonabelian Hodge graphs (untwisted case)},
J. Geom. Phys. \textbf{214} (2025), 105512.
\href{https://doi.org/10.1016/j.geomphys.2025.105512}{doi:10.1016/j.geomphys.2025.105512}.

\bibitem{BoalchDoucotRembado2025}
P.~Boalch, J.~Dou\c cot and G.~Rembado,
\emph{Twisted local wild mapping class groups: configuration spaces, fission trees and complex
braids},
Publ. Res. Inst. Math. Sci. \textbf{61} (2025), no.~3, 391--445.
\href{https://doi.org/10.4171/PRIMS/61-3-4}{doi:10.4171/PRIMS/61-3-4}.

\bibitem{BoalchYamakawa2020}
P.~Boalch and D.~Yamakawa,
\emph{Diagrams for nonabelian Hodge spaces on the affine line},
C. R. Math. \textbf{358} (2020), no.~1, 59--65.
\href{https://doi.org/10.5802/crmath.11}{doi:10.5802/crmath.11}.

\bibitem{BoalchYamakawa2015Twisted}
P.~Boalch and D.~Yamakawa,
\emph{Twisted wild character varieties},
arXiv:1512.08091 (2015).
\href{https://arxiv.org/abs/1512.08091}{arXiv:1512.08091}.

\bibitem{Doucot2021Diagrams}
J.~Dou\c cot,
\emph{Diagrams and irregular connections on the Riemann sphere},
arXiv:2107.02516 (2021).
\href{https://arxiv.org/abs/2107.02516}{arXiv:2107.02516}.

\bibitem{Doucot2025Graphs}
J.~Dou\c cot,
\emph{New nonabelian Hodge graphs from twisted irregular connections},
arXiv:2509.24861 (2025), revised 2026.
\href{https://arxiv.org/abs/2509.24861}{arXiv:2509.24861}.

\bibitem{Doucot2026Basic}
J.~Dou\c cot,
\emph{Basic representations of genus zero nonabelian Hodge spaces},
SIGMA \textbf{22} (2026), 027, 43 pp.
\href{https://doi.org/10.3842/SIGMA.2026.027}{doi:10.3842/SIGMA.2026.027}.

\bibitem{Hiroe2017Ramified}
K.~Hiroe,
\emph{Ramified irregular singularities of meromorphic connections and plane curve singularities},
Josai Math. Monogr. \textbf{10} (2017), 161--192.
\href{https://libir.josai.ac.jp/il/user_contents/02/G0000284repository/pdf/JOS-13447777-1011.pdf}{publisher PDF}.

\bibitem{Hiroe2019Rigidity}
K.~Hiroe,
\emph{Index of rigidity of differential equations and Euler characteristic of their spectral
curves},
arXiv:1907.12846 (2019).
\href{https://arxiv.org/abs/1907.12846}{arXiv:1907.12846}.

\bibitem{MartinezMichonSanMartin1994}
S.~Mart\'{i}nez, G.~Michon and J.~San Mart\'{i}n,
\emph{Inverse of strictly ultrametric matrices are of Stieltjes type},
SIAM J. Matrix Anal. Appl. \textbf{15} (1994), no.~1, 98--106.
\href{https://doi.org/10.1137/S0895479891217011}{doi:10.1137/S0895479891217011}.

\bibitem{McDonaldEtAl1995}
J.~J.~McDonald, M.~Neumann, H.~Schneider and M.~J.~Tsatsomeros,
\emph{Inverse $M$-matrix inequalities and generalized ultrametric matrices},
Linear Algebra Appl. \textbf{220} (1995), 321--341.
\href{https://doi.org/10.1016/0024-3795(94)00077-Q}{doi:10.1016/0024-3795(94)00077-Q}.

\bibitem{NabbenVarga1995}
R.~Nabben and R.~S.~Varga,
\emph{Generalized ultrametric matrices---a class of inverse $M$-matrices},
Linear Algebra Appl. \textbf{220} (1995), 365--390.
\href{https://doi.org/10.1016/0024-3795(94)00086-S}{doi:10.1016/0024-3795(94)00086-S}.

\bibitem{Nabben2026Acyclic}
R.~Nabben,
\emph{Inverse acyclic $Z$-matrices, matrices of tree structure and shifted generalized
ultrametric matrices},
arXiv:2607.27221 (2026).
\href{https://arxiv.org/abs/2607.27221}{arXiv:2607.27221}.

\end{thebibliography}

\end{document}
